% ===================================================================
%  图表第 2 号 — 人体。— 课间休息。— 游戏。
%  Traduction 2026 d'apres texto/io, controlee sur texto/fr, texto/en
%  et texto/es. Consigne : voir texto/zh/10-tubiao-01.tex.
% ===================================================================

%%K t02-apar-1 apar
\VUsaut{4.0mm}
\VUtitre{15.0pt}{60}{图表第 2 号}
\VUsaut{2.0mm}
\VUtitre{11.4pt}{0}{人体。--- 课间休息。}
\VUtitre{11.4pt}{0}{游戏。}

%%K t02-tit-0 sub
\VUtitre{13.2pt}{0}{人体。}
\VUsaut{2.4mm}
\VUcentre{10.2pt}{0}{\textit{（一堂简单的博物课。）}}
\VUsaut{3.0mm}

%%K t02-01-1 p
1. --- 人体的主要部分是：\VUgras{头}、\VUgras{躯干}和\VUgras{四肢}。

%%K t02-tit-1 sub
\VUpk{10.2pt}{40}{头。}
\VUsaut{3.0mm}

%%K t02-02-1 p
2. --- 头分两部分：\VUgras{头骨}，里面装着\VUgras{脑}，外面覆着
\VUgras{头发}；还有\VUgras{脸}。

%%K t02-03-1 p
3. --- 我们的图上既看不见头骨，也看不见脑；可是脸上重要的部分看得非常
清楚。

%%K t02-04-1 p
4. --- 首先是\VUgras{额头}，它显出强健的才智和不停活动的心思。
\VUgras{太阳穴}很开阔。\VUgras{眼睛}有神，睁得大大的。浓密的\VUgras{眉毛}
和长长的\VUgras{睫毛}护着它。

%%K t02-05-1 p
5. --- 这里还有\VUgras{鼻子}和\VUgras{鼻孔}。鼻子供我们闻气味和呼吸。鼻子
两旁是\VUgras{脸颊}，再往外是\VUgras{耳朵}。脸的下半部由\VUgras{嘴}和
\VUgras{下巴}组成。

%%K t02-06-1 p
6. --- 这两处看得不大清楚，因为\VUgras{胡子}和\VUgras{络腮胡}把它们遮住了。
不过我们知道嘴里有两片\VUgras{嘴唇}、\VUgras{上颚}和\VUgras{下颚}连同它们
的\VUgras{牙齿}，还有\VUgras{舌头}。我们用嘴说话吃饭。我们用舌头尝食物的
味道。

%%K t02-07-1 p
7. --- 眼、耳、鼻、口连同手指，是五种感觉的器官：\VUgras{视觉}、
\VUgras{听觉}、\VUgras{嗅觉}、\VUgras{味觉}和\VUgras{触觉}。

%%K t02-08-1 p
8. --- 所以头是人体最有意思的部分之一。

%%K t02-tit-2 sub
\VUsaut{4.4mm}
\VUpk{10.2pt}{40}{躯干。}
\VUsaut{3.0mm}

%%K t02-09-1 p
9. --- \VUgras{脖子}把头连在躯干上。它包括\VUgras{喉咙}和\VUgras{后颈}。

%%K t02-10-1 p
10. --- 躯干里也装着许多要紧的器官。\VUgras{胸}里有\VUgras{肺}，我们用它
呼吸；还有\VUgras{心}，那是生命的中心。心一停跳，活的生物就死了。

%%K t02-11-1 p
11. --- \VUgras{肋骨}撑着胸。

%%K t02-12-1 p
12. --- 躯干别的部分是\VUgras{腰侧}、\VUgras{背}、\VUgras{肩}和\VUgras{腹}
（肚子）。我们侧身或仰身睡觉，用肩膀扛东西。

%%K t02-13-1 p
13. --- 腹里有\VUgras{胃}和\VUgras{肠}。它们供我们消化食物。

%%K t02-tit-3 sub
\VUsaut{4.4mm}
\VUpk{10.2pt}{40}{四肢。}
\VUsaut{3.0mm}

%%K t02-14-1 p
14. --- 我们有四条\VUgras{肢}：两条\VUgras{胳膊}和两条\VUgras{腿}。我们用
胳膊拿、握、举、收东西；用腿站、走、跑。

%%K t02-15-1 p
15. --- \VUgras{肩}把胳膊连在身上。一块很结实的\VUgras{肌肉}，就是二头肌，
使胳膊活动；\VUgras{肘}把胳膊连在\VUgras{前臂}上。\VUgras{手腕}是
\VUgras{手}的关节，手的末端是\VUgras{手指}和\VUgras{指甲}。手上看得出一条
\VUgras{静脉}（还有一条动脉），\VUgras{血}在里面流。

%%K t02-16-1 p
16. --- 腿的各部分是：\VUgras{大腿}、\VUgras{膝}和\VUgras{腘窝}、
\VUgras{小腿肚}，它的\VUgras{肉}覆着\VUgras{皮}，还有\VUgras{脚踝}和
\VUgras{脚}。腿的\VUgras{骨头}很粗很壮。脚的末端是\VUgras{脚趾}。别的部分
是\VUgras{脚掌}和\VUgras{脚跟}。

%%K t02-17-1 p
17. --- 以上就是人体差不多完整的描述。

%%K t02-17-2 p
\textit{附注。} --- \textit{借着「我的……疼（头，等等）」这个说法，老师
可以从} \VUgras{身体健康}\textit{好坏的角度，把这一整套词汇再过一遍。}

%%K t02-tit-4 sub
\VUsaut{5.0mm}
\VUpk{10.2pt}{40}{体操。}
\VUsaut{2.4mm}
\VUcentre{10.2pt}{0}{\textit{（由一个学生讲述。）}}
\VUsaut{3.0mm}

%%K t02-18-1 p
18. --- 要身手灵便、结实、健康，就得练腿练胳膊。所以我们做游戏，也做
\VUgras{体操}。

%%K t02-19-1 p
19. --- 平日这两样都在学校的院子里做（见图表第 1 号）。可是星期四和星期日
我们到一片大草地上去，那里有地方跑，也有地方玩。

%%K t02-20-1 p
20. --- 老师和父母有时跟我们一起去；不过指挥这些练习的是
\VUgras{体操老师} \textsuperscript{(12)}。

%%K t02-21-1 p
21. --- 草地上，入口左边立着\VUgras{体操器械} \textsuperscript{(1)}：我们爬
\VUgras{梯子} \textsuperscript{(2)}，在\VUgras{吊环} \textsuperscript{(3)} 和
\VUgras{秋千架} \textsuperscript{(4)} 上倒着摆荡，两手交替攀到\VUgras{绳梯}
\textsuperscript{(5)} 或\VUgras{打结的绳子} \textsuperscript{(6)} 的顶上。

%%K t02-22-1 p
22. --- 同时我们的伙伴们跳过\VUgras{木马} \textsuperscript{(7)} 或
\VUgras{双杠} \textsuperscript{(11)}。另一些人在\VUgras{拳击} \textsuperscript{(8)}
或者戴着面罩拿着\VUgras{花剑}\VUgras{击剑} \textsuperscript{(9)}。还有一些人在
举\VUgras{哑铃} \textsuperscript{(10)}。

%%K t02-tit-5 sub
\VUsaut{4.6mm}
\VUpk{10.2pt}{40}{游戏。}
\VUsaut{3.0mm}

%%K t02-23-1 p
23. --- 体操的人四周，男孩女孩在玩各种\VUgras{游戏}。女孩们玩她们的
\VUgras{铁环} \textsuperscript{(14)}；她们\VUgras{跳绳} \textsuperscript{(13)}，或者
邀哥哥和表姐妹来一局\VUgras{槌球} \textsuperscript{(15)}。对手正以为能轻轻松松
过门的时候，她们用\VUgras{木槌}把他的\VUgras{球}打开，别提多得意了！

%%K t02-24-1 p
24. --- 男孩们更爱使劲的游戏。他们乐意打\VUgras{九柱} \textsuperscript{(16)}，
用\VUgras{球} \textsuperscript{(17)} 干净利落地把柱子撞倒。他们也不嫌
\VUgras{陀螺} \textsuperscript{(18)} 和\VUgras{接球棒} \textsuperscript{(19)}，甚至
\VUgras{投蛤蟆} \textsuperscript{(20)}。不过他们最爱的游戏是\VUgras{弹珠}
\textsuperscript{(21)} 和\VUgras{手球} \textsuperscript{(22)}，因为\VUgras{暖房}
\textsuperscript{(26)} 的墙正好拿来弹那只轻球。

%%K t02-25-1 p
25. --- 小约翰踩着他的\VUgras{高跷} \textsuperscript{(24)}，很有本事，站得
稳稳的。他旁边有几个朋友在那座暖房里和\VUgras{篱笆}后面玩\VUgras{捉迷藏}
\textsuperscript{(25)}。另一些人更爱\VUgras{猜手掌} \textsuperscript{(28)} 或者
\VUgras{毽球} \textsuperscript{(29)}，那毽球从一只\VUgras{球拍}飞到另一只。

%%K t02-noto-1 noto
\VUnotes{143.2mm}{%
(*) 儿童游戏往往只有各国自己的名字。我们尽力用伊多语给它们命名：有时取
它们最显著的动作，有时在不至于叫人误会时采用某一国的名字。例如：
\VUgras{木桶戏}（德国和法国的）就是\VUgras{投蛤蟆} \textsuperscript{(20)}
（英国的），玩的人要把圆片投进那只张着嘴、蹲在带孔箱子（木桶）上的
金属蛤蟆口里。\VUgras{越肩跳} \textsuperscript{(30)} 和\VUgras{跳背}只有一点
不同：跳过头顶时，玩的人两手撑在同伴的肩上，而在「\VUgras{跳背}」里只撑
在他的背上。——又或者，一些人弯下腰，另一些人两手撑着他们的肩跳过他们的
背；前者要顶住冲力不塌下去，后者要跳过去不失去平衡（见挂图本身）。%
}

%%K t02-26-1 p
26. --- 草地当中有两棵树。其中一棵树下，七八个男孩玩起了\VUgras{越肩跳}
\textsuperscript{(30)}(*)。这游戏不是每个国家都有；可是人人都知道什么是
\VUgras{跳背}，孩子们这会儿没有玩它。

%%K t02-tit-6 sub
\VUsaut{5.0mm}
\VUpk{10.2pt}{40}{露天的游戏。}
\VUsaut{3.4mm}

%%K t02-27-1 p
27. --- \VUgras{捉瞎子} \textsuperscript{(31)} 也很好玩；女孩男孩轮流把
\VUgras{眼罩}蒙在眼上，去捉那些笨手笨脚、让人逮住的人。

%%K t02-28-1 p
28. --- 草地当中有一种很法国的游戏，虽说粗鲁了些：这就是
\VUgras{熊} \textsuperscript{(32)}，比安安静静的\VUgras{风筝} \textsuperscript{(33)}
吵闹得多，也比英国的\VUgras{网球} \textsuperscript{(34)} 吵闹得多。

%%K t02-29-1 p
29. --- 最后这项运动是大孩子们的乐趣。我们的老师自己也乐意来上一场。它要
很大的本事和灵活劲儿，我很喜欢它。\VUgras{秋千} \textsuperscript{(35)} 我留给
女孩们。

%%K t02-30-1 p
30. --- 另一种英国游戏我也不喜欢，它很可能变得危险。

%%K t02-30-1b p
我说的是\VUgras{足球} \textsuperscript{(36)}，那是我哥哥的乐事。我更爱我们的
老游戏\VUgras{夺俘} \textsuperscript{(38)}，一样有益身体，却粗野得多。

%%K t02-tit-7 sub
\VUsaut{4.6mm}
\VUpk{10.2pt}{40}{骑车和球类游戏。}
\VUsaut{3.0mm}

%%K t02-31-1 p
31. --- 我也乐意骑\VUgras{自行车} \textsuperscript{(37)} 绕着草地跑。

%%K t02-32-1 p
32. --- 明年我要学\VUgras{骑马} \textsuperscript{(39)}。马迈着漂亮的步子或小跑
起来，又纵着蹄子跨过障碍，多好看啊！

%%K t02-33-1 p
33. --- 夏天我们学\VUgras{游泳} \textsuperscript{(40)}。我们在河里清凉的水中
痛痛快快地扎猛子、洗澡。有时到末了，我们放一只纸\VUgras{气球}
\textsuperscript{(41)}，它一装满热气就庄严地升上天空。

%%K t02-34-1 p
34. --- 还有许多别的游戏，可是我数也数不完，从这段短短的话里也看得出，
在我们那片好草地上，星期四一点儿也不闷。

%%K t02-34-2 p
N.~B. --- \textit{老师很容易把这一章补足，只要把比较要紧的游戏一个一个
讲得更细些。}
\VUsaut{6.0mm}
\VUfilet{22mm}
